\pdfoutput=1
\documentclass[11pt]{article}
\usepackage{acl2023}
\usepackage{times}
\usepackage{latexsym}
\usepackage[T1]{fontenc}
\usepackage[utf8]{inputenc}
\usepackage{microtype}
\usepackage{inconsolata}
\usepackage{booktabs}
\usepackage{array}
\usepackage{xcolor}

\title{Named Entity Recognition for Information Security and Cyber Threat
Intelligence: A Survey}

\author{Daniel Dye \\
  Bellini College of Artificial Intelligence, Cybersecurity, and Computing \\
  University of South Florida \\
  \texttt{ddye@usf.edu}}

\begin{document}
\sloppy
% displaCy-style inline entity annotation: highlighted span + small-caps label
\newcommand{\ent}[3]{\colorbox{#1}{\strut #2\,{\tiny\textsc{#3}}}}

\maketitle

\begin{abstract}

Cyber threat intelligence is produced as natural language and consumed as structured data; named entity recognition and its downstream tasks — entity linking, relation extraction, TTP mapping, and knowledge-graph construction — form the extraction bottleneck between the two. We survey this pipeline across 41 works spanning classic sequence labeling, weak and distant supervision, zero-shot span matching (the GLiNER family), LLM-based generative extraction, data-centric methods, and the 2024-2026 wave of LLM-driven knowledge-graph construction, together with the taxonomies and ontologies (STIX, MITRE ATT\&CK, UCO) that serve as extraction type systems and the benchmarks that evaluate the results. Two findings organize the synthesis. First, fragmentation recurs at every level — annotation schemas, entity names, KG ontologies, vendor coverage — with empirical costs now quantified (cross-dataset F1 collapse; thousands of unreconciled actor aliases) and an economic root: convergence requires sharing the proprietary telemetry that justifies vendors' analytic claims. Second, the field's methodological diversity is largely a single response to one condition — expert annotation too expensive to scale — pursued by substituting existing structure for labels: knowledge bases as distant supervision, type descriptions as zero-shot prompts, canonical identifiers as verifiable rewards, and LLM-generated text as synthetic corpora. We catalog reproducible artifacts for every method, identify the benchmark gaps that matter operationally (consistency, calibration, relation extraction, label efficiency), and chart open problems from adversarial robustness of ingestion to schema-agnostic graph construction.

\end{abstract}
\section{Introduction}

Cyber threat intelligence (CTI) is produced and consumed primarily as natural language. Government advisories, vendor blog posts, incident reports, and malware analyses describe threat actors, malware families, exploited vulnerabilities, and adversary techniques in prose — and the security tooling that must act on this knowledge (detection rules, threat feeds, knowledge graphs, attack simulations) requires it as structured data. Named entity recognition sits at this bottleneck: it is the extraction step that turns a sentence like "Volt Typhoon used living-off-the-land binaries to evade detection" into typed, linkable objects:

\begin{center}
\small
\setlength{\fboxsep}{2.5pt}
\ent{orange!25}{Volt Typhoon}{threat actor group} used
\ent{cyan!25}{living-off-the-land binaries}{attack technique name}
to \ent{green!25}{evade detection}{tactic}.
\end{center}

The practitioner pipeline is consistent across organizations: collection (scraping advisories and reports), extraction (NER and related tasks), structuring (mapping to standards such as STIX 2.1 and MITRE ATT\&CK), and use (feeds, hunting, correlation). Every downstream stage inherits the quality of the extraction stage.

CTI NER differs from general-domain NER in ways that make it a genuinely distinct research problem rather than a domain transfer exercise. The entity inventory is unusual: alongside conventional organizations and locations, CTI text contains malware families, threat-actor groups with contested naming conventions, file hashes, registry keys, CVE identifiers, YARA rule names, and adversary techniques that may be mentioned implicitly rather than named \citep{lange2024annoctr}. The vocabulary is adversarial and fast-moving — new malware names appear weekly, guaranteeing persistent out-of-vocabulary pressure. And the supervision economics are poor: annotation requires scarce domain expertise, and most existing corpora are small and ambiguously licensed \citep{lange2024annoctr}.

This survey synthesizes the literature on NER for information security and CTI, with particular attention to developments from 2024–2026 that have reshaped the field: the collision between dataset fragmentation and harmonization efforts \citep{jalocha2025label,echchammakhy2025cyberner}, the arrival of compact zero-shot span-matching architectures~\citep{zaratiana2023gliner} that sidestep fixed label schemas entirely, LLM-based extraction and the growing evidence of its reliability limits, and the maturation of TTP mapping — assigning text to MITRE ATT\&CK techniques — as an adjacent extraction task with its own methods and benchmarks.

A thread we develop throughout: many of the field's apparently separate research directions — label unification, schema harmonization, distant supervision, synthetic data generation, zero-shot architectures, instruction refinement, verifiable-reward training — are best understood as responses to a single underlying condition: \emph{expert CTI annotation is too expensive to produce at the scale modern models require.} The field has repeatedly developed methods to avoid requiring it, most often by substituting structure that already exists (knowledge bases, type descriptions, canonical identifiers, the standards themselves) for labels that do not.

\section{Background}

\subsection{Task definitions}

Span-level NER is the core task: given text and a set of entity types, return typed spans — classically formulated as BIO/BIOES sequence labeling (the convention in CyNER, DNRTI, APTNER; BIOES-to-BIO conversion is itself a harmonization step, \citealp{echchammakhy2025cyberner}), and more recently as span-type matching where candidate spans are scored against type representations (the GLiNER family, §5.5), which natively supports nested and overlapping entities that tagging schemes cannot express (GLiNER2-PII's full\_name containing first\_name). Entity linking extends NER by resolving mentions to knowledge-base identities — Wikipedia/Wikidata for general entities, ATT\&CK objects for techniques (\citealp{lange2024annoctr,hanks2022recognizing}) — and becomes load-bearing in CTI because surface forms are unreliable designators (§2.3's naming problem). Relation extraction types entity pairs, jointly (GLiNER-Relex) or in pipeline, and is the step that turns extraction output into graph triples (§7).

Two CTI-specific complications cut across these definitions. First, the mention/concept distinction: AnnoCTR shows adversary techniques are frequently \emph{implicit} — described periphrastically with no taggable name span ("send phishing e-mails with malicious attachments" for T1566.001\footnote{AnnoCTR's text prints T0865, the ICS-matrix counterpart; we give the Enterprise ATT\&CK identifier for Spearphishing Attachment.}) — which is why TTP mapping (§6) is better formulated as classification or linking at document, sentence, or span granularity (the rcATT-to-TRAM-to-AnnoCTR progression) than as NER. Second, span-length conventions diverge by purpose: IOC-centric corpora annotate noun phrases, while attribution-centric AttackER deliberately annotates phrase- and sentence-length spans (COURSE\_OF\_ACTION, IMPACT) — a design choice, not noise, but one that breaks cross-corpus comparability (§3).

\subsection{Method eras}

The methodological history of cybersecurity NER recapitulates general NLP, with a persistent domain-specific twist: rule-based extraction never went away. Early systems combined gazetteers and regular expressions — which remain genuinely effective for structurally regular entities like CVE identifiers, file hashes, and IP addresses — with feature-engineered statistical models, where conditional random fields became the sequence-labeling standard~\citep{echchammakhy2025cyberner}. The deep learning era replaced feature engineering with BiLSTM-CRF architectures augmented by character-level features and domain-specific word embeddings. Transformer fine-tuning then became dominant, accompanied by domain-pretrained encoders: SecureBERT (trained on 2.2M security documents), CySecBERT, and DarkBERT (dark-web corpora), typically with CRF decoding heads~\citep{echchammakhy2025cyberner}.

Two newer paradigms now compete with — and compose with — this lineage. Zero-shot span-matching models (GLiNER family, \citealp{zaratiana2023gliner}) express entity types as natural-language prompts matched against span representations, removing the fixed-schema constraint that made every new CTI label set a retraining problem. Generative LLM extraction offers the same schema flexibility through instruction-following, at substantially higher inference cost and with reliability concerns specific to CTI that we examine in §5.6 and §8.

In production CTI systems, these families are typically \emph{ensembled} rather than substituted: regex and KB lookup provide high-precision extraction of structured indicators, while learned models contribute contextual entities (victim organizations, sectors, implicitly described techniques) that rules cannot reach.

\subsection{CTI taxonomies, standards, and ontologies}

Every CTI extraction system commits, implicitly or explicitly, to some taxonomy as its type system — so the state of CTI knowledge representation is not background to this survey but a determinant of everything in it: taxonomy choice fixes what counts as an entity (§3), which ontology a knowledge graph populates (§7), and whether two systems' outputs are even comparable (§8). The operational landscape divides labor roughly as follows. STIX 2.1 provides the exchange format and object model — domain objects (Threat-Actor, Malware, Tool, Campaign, Intrusion-Set) and cyber-observables (IPv4-Addr, File, Domain-Name) — and has become the de facto NER type system: CyberNER harmonizes onto it~\citep{echchammakhy2025cyberner}, AttackER derives its schema from it~\citep{deka2024attacker}, and this corpus's own annotation pipeline extracts from its bundles. MITRE ATT\&CK supplies the behavioral taxonomy (14 tactics, $\sim$193 techniques, $\sim$401 sub-techniques) that the TTP-mapping literature of §6 targets. Around these two sit the specialized inventories Mavroeidis and Bromander catalog~\citep{mavroeidis2021cyber}: CAPEC for attack patterns, CWE/CVE for weaknesses and vulnerabilities, MAEC for malware characterization, VERIS for incident metrics, MISP's galaxy taxonomies for sharing communities, and the formal-ontology lineage (Unified Cybersecurity Ontology, CASE) that attempts to give all of this actual semantics.

The standing verdict on that landscape comes from Mavroeidis and Bromander's evaluation, and its conclusions have proven durable: scored against a model of what CTI must express (the who, what, why, where, when, and how of threats, plus indicators and courses of action), existing taxonomies, standards, and ontologies are "not thoroughly designed, non-interoperable, ambiguous, and lacking proper semantics and axioms for reasoning." Concept ambiguity blocks ontology integration; prose-heavy fields undermine queryability; missing relationships between concepts prevent explainable views. First published in 2017 and revised in 2023 with conclusions intact, the evaluation describes problems the field has since confirmed rather than resolved — the label-schema failures of §3 and the per-system KG ontologies of §7 are downstream symptoms. Recent formal work sharpens the critique at the foundations: a UFO-based analysis (MAECO-Lite, \citealp{adams2026maecolite}) shows that MAEC and STIX conflate enduring artifacts with runtime events — the malware sample versus what it did when executed — and that this representational sloppiness measurably degrades machine reasoning (description-logic concept learning improves significantly on a foundationally clean modular replacement). The type systems NER extraction targets are, in other words, themselves ontologically confused.

Fragmentation also operates below the schema level, at the instances. The "Hesperus is Phosphorus" study~\citep{roa2025hesperus} maps threat-actor naming at scale — 13,371 reports, 17 vendor taxonomies, 3,287 actor names across 15 sources — and finds aliases concentrated on a small set of prominent actors, proliferating over time, with existing inter-vendor mappings containing outright errors inflated by temporary activity-cluster names, shared infrastructure, and overlapping operations. The structural insight matters more than the counts: vendors cannot converge on names because agreeing that one vendor's APT29 is another's Midnight Blizzard requires sharing the private telemetry behind the attribution. Naming fragmentation is thus not carelessness but an equilibrium of the CTI market — the same economics that keep corpora small and unlicensed (§4) keep the namespace fragmented. For extraction, the consequence is concrete: an ACTOR span is only as useful as its normalization, and cross-source entity resolution inherits a problem the industry itself has not solved.

The current SOTA direction treats ontologies not as passive type systems but as active constraints on extraction. OntoLogX~\citep{cotti2025enabling} wires a domain ontology and SHACL shape validation directly into LLM-based extraction from security logs: the ontology defines the output structure, the constraints enforce semantic validity over the resulting graph, and accuracy improves over prompt-only baselines. The pattern — formal semantics as guardrails on neural extraction — anticipates both the LLM-reliability mitigation question of §5.6 and the ingestion-validation stage of KG construction (§7), and the 2026 ontology-discovery work covered there~\citep{kim2026schemaagnostic} closes the loop by asking whether fixed schemas should be replaced by schemas learned from the data.

\section{The Schema Problem: Label Fragmentation and Harmonization}

If one finding from the recent literature should reframe how researchers approach CTI NER, it is this: the field's datasets do not merely use different labels — they encode \emph{incompatible conceptualizations} of what the entities are, and this incompatibility has severe, measurable consequences. Ja\l{}ocha et al.~\citep{jalocha2025label} unified four prominent corpora (APTNER, CyNER, DNRTI, AttackER) onto four coarse labels and measured pairwise cross-dataset performance with matched architectures. Off-diagonal span-F1 ranged from 0.01 to 0.31 against in-dataset baselines of 0.23–0.41. More strikingly, training on the union of all four unified datasets \emph{reduced} performance on two of them relative to single-dataset training — the noise introduced by unification dominated the benefit of quadrupled training data.

The error analysis identifies why. The same token receives contradictory gold types across corpora: 'sample' maps (post-unification) to Malware in DNRTI but O in APTNER; 'backdoor' and even the specific malware name 'Dridex' are Malware in APTNER but System in DNRTI; operating systems are System in APTNER and CyNER but O in DNRTI. These are not annotation errors — each is locally consistent with its corpus's guidelines. They are definition discrepancies: each small dataset crystallized its annotator's local ontology, and no dataset was large or authoritative enough to become a standard. Span-length conventions diverge too (AttackER favors long spans where others annotate noun phrases), and Jensen-Shannon divergence between corpora's word/POS distributions correlates strongly (Pearson $\approx$ $-$0.74) with cross-dataset degradation.

CyberNER~\citep{echchammakhy2025cyberner} demonstrates that the failure is not inherent to dataset combination but to \emph{ad-hoc} combination. Harmonizing the same four corpora onto STIX 2.1 — consolidating synonymous tags (APT, HackOrg, THREAT\_ACTOR $\rightarrow$ Threat-Actor), differentiating conflated ones (Malware vs Tool vs Software), mapping technical indicators to STIX cyber-observable types, and explicitly excluding unmappable tags — produces a 610k-token corpus on which models gain roughly 30\% relative F1 over naive concatenation. The choice of STIX as the target taxonomy matters beyond the numbers: extraction output becomes directly consumable by STIX-speaking CTI infrastructure, and the ontology's defined relationships open a path to relation extraction and knowledge-graph construction.

Harmonization is nonetheless lossy. CyberNER's exclusion strategy discards source tags that resist STIX mapping; fine-grained distinctions (APTNER's separate MD5/SHA1/SHA2 types collapse into File) survive only as recoverable structure, and the explicit/implicit mention distinction that AnnoCTR shows matters for technique extraction has no harmonized representation. An alternative line keeps original label sets and shares representations instead: Ja\l{}ocha et al.'s multi-head model (per-dataset output heads over shared embeddings) matches unified-label performance on most datasets while preserving source-schema specificity — suggesting that the field's fragmentation may be better managed than eliminated.

\section{Datasets and Corpora}

The corpus landscape divides into a fragmented first generation and a consolidating second one. The first generation — MalwareTextDB, CTI-Reports, rcATT, CyEnts, CySecED, CASIE, and the four corpora discussed in §3 — is characterized by small document counts (tens to hundreds), custom schemas, and, critically, absent or unclear licensing: AnnoCTR's survey of the landscape \citep[Table 1]{lange2024annoctr} finds that nearly every prior dataset lacks a usable license, blocking replication, redistribution, and any commercial application. The second generation addresses these defects directly: AnnoCTR secured CC-BY-SA release by having CTI vendors donate 400 reports, and CyberNER manufactured scale (610k tokens) by harmonizing existing corpora rather than annotating anew.

\begin{table*}[t]
\centering
\footnotesize
\setlength{\tabcolsep}{4pt}
\begin{tabular}{>{\raggedright\arraybackslash}p{3.0cm}>{\raggedright\arraybackslash}p{2.3cm}>{\raggedright\arraybackslash}p{2.4cm}>{\raggedright\arraybackslash}p{2.1cm}>{\raggedright\arraybackslash}p{1.5cm}>{\raggedright\arraybackslash}p{2.7cm}}
\toprule
Corpus & Size & Entity types & Source texts & License & Notes \\
\midrule
CyNER (2022) & 107k tok / 7.6k ents & 5 & 60 TI reports & open & coarse labels; BIO \\
DNRTI (2020) & 176k tok / 36.9k ents & 12-13 & TI reports, gov, GitHub & unclear & BIO \\
APTNER (2022) & 260k tok / 39.7k ents & 21-23 & vendor reports & unclear & BIOES; fine-grained (MD5, VULID) \\
AttackER~\citep{deka2024attacker} & 67-79k tok / 31.1k ents (7,026 instances) & 18 (STIX-derived + 2) & 217 attribution reports & CC-BY-4.0 (Zenodo) & long spans by design \\
AnnoCTR~\citep{lange2024annoctr} & 400 docs / 120 cyber-annotated / 20,855 ann.\ (13,244 cyber layer) & GNE + CyNE + TACTIC\slash TECHNIQUE + TIMEX, CodeSnippet & 5 vendor blogs, donated & CC-BY-SA 4.0 & explicit/implicit; ATT\&CK + Wikipedia links \\
CyberNER~\citep{echchammakhy2025cyberner} & 610k tok / 103k ent tokens & 21 (STIX 2.1) & union of first four & public & harmonized; 16.93\% density \\
AZERG~\citep{lekssays2025text} & 141 reports / 4,011 ents / 2,075 relations & STIX entities + relations & full threat reports & Apache-2.0 (HF) & rare relationship annotations \\
CTIfecta (via \citealp{haque2026single}) & 106 reports, 3 campaigns & ATT\&CK techniques & multi-vendor campaign reports & public & campaign-level, multi-report \\
Legacy (MalwareTextDB, CTI-Reports, rcATT, CyEnts, CySecED, CASIE, TRAM) & small & various & various & mostly none & landscape per AnnoCTR Table 1 \\
\bottomrule
\end{tabular}
\caption{CTI NER corpora. First-generation corpora are characterized by
small document counts, custom schemas, and absent or unclear licensing;
second-generation corpora optimize legal reusability, interoperability, and
evaluation realism.}
\label{tab:corpora}
\end{table*}


(AttackER token counts differ between Ja\l{}ocha et al. and CyberNER — 67k vs 79k — likely a versioning difference; the table reports both figures.)

Read together, the table records a generational shift in what corpus builders optimize for. The first generation optimized annotation density on whatever text was at hand, producing the schema and licensing chaos sections 3 and 9 document. The second optimizes legal reusability (AnnoCTR's donated-text model), interoperability (CyberNER's STIX anchoring, AttackER's STIX derivation), relationship coverage (AZERG), and evaluation realism (CTIfecta's multi-report campaign structure). A third-generation corpus would add what none yet has: deliberate source diversity against vendor bias, ingestion verification against adversarial text, and span annotations designed for both NER and linking evaluation.

Annotation methodology remains the landscape's weakest element. AnnoCTR split annotation by expertise, using a two-person trained team for its general layer and a single domain expert for its cybersecurity layer (INCEpTION), and its rough second-annotator study quantifies why the field so rarely reports agreement: a second security-literate professional, working without an extended shared training phase, reached span F1 of only 0.31 on CONCEPT, 0.52 on GROUP, and 0.67 on MALWARE, agreed at 0.54 on document-level technique sets, and marked no TOOL spans at all \citep{lange2024annoctr}. APTNER/DNRTI used Brat with overlapping author teams (one explanation Ja\l{}ocha et al. offer for their unusual mutual transferability); AttackER mixed manual and semi-automatic passes; inter-annotator agreement is otherwise rarely reported anywhere. The economics explain it: CTI annotation needs scarce dual expertise, which is the root condition this survey's thesis traces through every method family. Two structural source problems compound the corpus problems: vendor coverage barely overlaps and carries geographic/sectoral bias (the CTI echo chamber, \citealp{suarezroman2026cti} — 16,096 reports analyzed; a few core providers give most coverage, additional sources yield diminishing returns), so corpus source selection is a representativeness decision, not a convenience one (the skew operates even within a single corpus: ZScaler's 27 documents contribute 54\% of AnnoCTR's implicit technique mentions); and ingestion is adversarially contestable — generated fake-intelligence text passes both specialized classifiers (97\% FPR) and GPT-4o filtering (75\% FPR)~\citep{shafee2025false}, meaning corpus pipelines that scrape open sources need verification stages the current generation lacks.

\section{Methods}

Organizing the methods literature by \emph{supervision source} and \emph{type-binding mechanism} (rather than chronology) yields seven families; the first six are model families, the seventh operates on data and composes with any of them.

\subsection{Rule- and dictionary-based extraction}

Rule-based extraction has remained in continuous use in CTI NER, for principled reasons. A substantial fraction of CTI entities are \emph{structurally regular} — CVE identifiers, hashes, IP addresses, ATT\&CK T-codes, registry keys — where regular expressions deliver near-perfect precision at zero training cost with full auditability (\citealp{echchammakhy2025cyberner}; the practical lineage runs through \citet{hanks2022recognizing} and CyNER's own indicator handling, which Ja\l{}ocha et al. note was deliberately excluded from learned labels because regex suffices). Gazetteer and knowledge-base lookup extends the same logic to proper-noun inventories: MITRE ATT\&CK's group/tool/malware lists provide thousands of canonical names and aliases. The family's limits are equally principled: no contextual disambiguation (short common-word aliases generate false positives), no coverage of unlisted entities, and maintenance burden as vocabularies evolve. In modern systems rules are an ensemble layer, not a competitor — and via distant supervision (§5.4) they became a mechanism for generating training data; the bi-encoder retrieval formulation (§5.5) is their latest reinvention, replacing string match with semantic match against the same inventories.

\subsection{Classical statistical sequence labeling}

The feature-engineered era — HMMs and SVMs giving way to CRFs as the sequence-labeling standard~\citep{echchammakhy2025cyberner} — matters to this survey mostly for what it established: the BIO formulation, the per-token feature mindset that domain embeddings later replaced, and CRF decoding heads that survived into the neural era. Its lasting lesson is that transition modeling (label dependencies) adds value over independent token decisions, which is why transformer+CRF stacks remained competitive long after contextual embeddings arrived.

\subsection{Supervised neural sequence labeling}

BiLSTM-CRF with character features and security-corpus word embeddings became the de facto standard (the architecture of the original corpus papers' baselines, and the testbed of the cross-dataset study, \citealp{jalocha2025label}), displaced by transformer fine-tuning. The domain-adaptation move followed: SecureBERT (2.2M security documents), CySecBERT, and DarkBERT (dark-web corpora) consistently improve over general encoders, with transformer+CRF reporting state-of-the-art supervised results on the standard corpora~\citep{echchammakhy2025cyberner}. Two qualifications matter. First, supervised SOTA is corpus-bound: the same models that excel in-dataset collapse cross-dataset (§3), so "supervised SOTA" claims are claims about a label distribution, not about CTI text. Second, the paradigm's annotation appetite is precisely what the rest of this section's families exist to avoid — and within-family evidence shows the appetite can be engineered down (LETNER's 10\%-of-annotations result, §5.5; TTPrompt's 1\% instruction refinement, §5.6).

AnnoCTR's baselines add a third qualification from a controlled setting \citep{lange2024annoctr}: the identical RoBERTa tagger scores 82.3 micro-F1 on the corpus's general entities but only 57.8 on its cybersecurity entities, a 24-point gap with architecture and training data held constant. In the same comparison, generic RoBERTa outperformed both SciBERT and CodeBERT, which the authors attribute to pretraining-data register (RoBERTa's news-heavy corpus resembles threat-report prose) rather than domain vocabulary — a reminder that "domain-pretrained" conflates at least two axes, register and terminology, and that proximity of register can matter more than proximity of subject.

A concrete instance of the domain-adaptation move going further than swapping the encoder: SSNER retrains SecureBERT's own masking objective rather than just fine-tuning on top of it. Standard MLM masks independent tokens, which under-represents multi-token spans — a structural problem for CTI, where entities are disproportionately multi-word (malware family names, technique descriptions, product-and-version strings). SSNER's Segment Masking Enhancement (SME) masks contiguous random-length segments during a threat-intel-corpus retraining pass instead, then adds a Semantic Collaborative Embedding (SCE) module that fuses a filtered set of semantically similar words per mention via self-attention to address polysemy (the same entity string denoting different types in different contexts, discussed further in §9). On two established CTI NER benchmarks, DNRTI and Bridges et al., the combination reaches F1 86.12 and 96.02 respectively, beating the prior best models NER4CTI (+0.78) and FIEBD (+1.46); an ablation isolates each component's contribution (removing SME costs 2.88–3.40 F1, removing SCE costs 1.60–2.90 F1), and per-entity results correlate with multi-token share exactly as the mechanism predicts — categories that are 70–72\% multi-token entities score F1 0.95–0.97, among the model's best. The result is a data point against treating "domain-pretrained backbone" as a single lever: \emph{how} a backbone is retrained (objective, not just corpus) is a separable axis from \emph{what} it is fine-tuned on afterward.

\subsection{Weak and distant supervision}

Distant supervision generates training labels by matching corpus text against external knowledge bases, trading annotation cost for label noise. BOND~\citep{liang2020bond} established both the canonical framework and the precise characterization of the noise: incomplete annotation (KB coverage gaps leave real entities labeled as non-entities) and inaccurate annotation (context-free matching mislabels ambiguous mentions — 'Liverpool' the city versus the football club). Its two-stage recipe — fine-tune a pre-trained encoder on distant labels with aggressive early stopping, then discard the distant labels and self-train on teacher-generated pseudo-labels — remains the backbone of the lineage that followed (SCDL, ATSEN, CENSOR).

The frontier of that lineage is calibration. CENSOR~\citep{si2023improving} demonstrates that the confidence thresholds all prior teacher-student methods use to select "reliable" pseudo-labels are built on a false premise: neural teachers are poorly calibrated, so high confidence does not imply correctness, and confidently wrong pseudo-labels propagate through self-training. Its remedies — uncertainty estimation via prediction variance across stochastic forward passes, selection requiring both high confidence and low uncertainty, and paired student networks that exchange small-loss labels so mislabeled samples can be corrected rather than merely discarded — yield consistent gains (average +1.87 F1 across five DS-NER benchmarks).

Two observations matter for CTI specifically. First, BOND's own data shows distant supervision quality is a function of gazetteer cleanliness: KB matching achieves F1 $\approx$ 70 in biomedical domains with curated ontologies but $\approx$ 36–60 in the open domain. CTI sits in the favorable regime — STIX bundles, MITRE ATT\&CK, and CVE registries are exactly the kind of curated, low-ambiguity vocabularies that make distant labels usable. Second, ensemble-agreement supervision — using the intersection of independent annotators (dictionary and zero-shot NER) as silver labels — is a distant-supervision variant whose noise model differs from KB matching, and the calibration critique applies to it directly: agreement is a confidence-style heuristic, and uncertainty-aware selection over the ensemble is an unexplored refinement.

The BOND lineage has since branched along a third axis besides calibration (CENSOR) and case-study application: ensembling. BOND-MoE trains K independently-initialized experts on disjoint document subsets (hard-EM assignment, with a Sinkhorn-scaled fairness constraint preventing all documents collapsing onto one expert) before self-training the ensemble as in BOND's second stage — noise reduction through model diversity rather than through better confidence estimates. The result across five general-domain DS-NER benchmarks is genuinely mixed: gains on three (CoNLL03 +2.19, WebPage +0.42, Wikigold +0.83 F1 over BOND) and small losses on two (Twitter $-$0.31, OntoNotes $-$0.09), with the paper's own ablation showing the fairness-assignment constraint contributes more than the expert-ensembling idea in isolation. This is worth stating explicitly: ensembling-for-denoising is a live but not yet clearly superior alternative to CENSOR's uncertainty-based route, and the two have not been compared head-to-head.

A CTI-scoped instantiation of the BOND recipe exists, though its scope is narrower than the general framing above suggests. Zhang et al. apply distant supervision plus self-training denoising to exactly two entity types — Vendor and Product — extracted from 934 NIST NVD CVE descriptions, using a Reverse-Maximum-Matching-plus-POS-constraint algorithm (RMP) against two hand-built dictionaries for the distant labels, and a confidence-gated teacher-student loop (IST) with explicit weight-update constraints for denoising. Averaged across five backbones, IST improves F1 over raw distant labels by 3.5 points (Vendor) and 3.35 points (Product); the pretrained-transformer backbones (RoBERTa, RoBERTa-wwm-RDCNN-CRF) generalize best under self-training, consistent with BOND's own original finding that PLM-based methods outperform LSTM/CNN architectures on distant labels. The paper's own stated limitation — two hand-built lexicons cap label-type coverage — is precisely the scalability question a STIX/MITRE-anchored dictionary of the kind used in this repo's pipeline is positioned to answer at a scope closer to a full CTI schema than Vendor/Product alone.

\subsection{Zero-shot and few-shot span matching}

GLiNER~\citep{zaratiana2023gliner} reframes open-type NER as similarity matching rather than generation: entity types, expressed in natural language, are encoded jointly with the input text by a small bidirectional encoder; span representations are scored against type embeddings with a sigmoid dot-product. The architectural consequences matter for CTI deployment. All span-type pairs are scored in parallel (no sequential decoding); the model is compact (50M–300M parameters versus 7B–13B for fine-tuned LLM extractors); and new entity types are introduced by writing a phrase, not by retraining. On out-of-domain zero-shot benchmarks, GLiNER-L (0.3B) outperforms ChatGPT (60.9 vs 47.5 average F1) and fine-tuned 7B–13B instruction models.

For CTI specifically, the zero-shot framing is a direct answer to the schema problem of §3: if no fixed label set is baked into the model, fragmented corpus schemas become prompt-engineering rather than retraining problems. The trade-offs — sensitivity to type-name phrasing, span-length ceilings (K=12 tokens), and weaker performance on structurally regular entities like hashes that rules capture trivially — motivate the ensemble deployments noted in §2.2.

The architecture has since become an ecosystem, and its trajectory tracks three pressures relevant to CTI deployment. The first is task unification: GLiNER2~\citep{zaratiana2025gliner2} folds NER, text classification, and hierarchical structured extraction into one 205M schema-driven model, where extraction tasks are declared as input schemas rather than baked into architecture — and, critically for CTI, entity types carry natural-language \emph{descriptions}, making ATT\&CK technique definitions first-class model inputs rather than augmentation tricks. The second is label-set scale: joint encoding couples inference cost to label count, which breaks at CTI-sized vocabularies (hundreds of ATT\&CK techniques, thousands of malware and actor names). The bi-encoder variant~\citep{stepanov2026millionlabel} decouples label and context encoding so label embeddings are pre-computed and cached — constant-time in label count, up to 130x throughput at 1024 labels, while \emph{raising} zero-shot CrossNER to 61.5 micro-F1 — and its GLiNKER linking framework turns gazetteer lookup into semantic retrieval against knowledge bases. The third is domain adaptation under data scarcity, where the PII variant~\citep{zaratiana2026gliner2pii} matters for what it proves rather than its domain: trained entirely on constraint-driven synthetic data (4,910 examples), it beats four established detectors on natural-text benchmarks. Schema-conditioned encoders have also reached multi-aspect safety classification (GLiGuard, \citealp{zaratiana2026gliguard}), confirming the paradigm's reach beyond extraction.

The span-matching family has already drawn CTI-specific architectural critique. LETNER (Wang et al., ESWA 2026) argues that GLiNER's span representations — built by concatenation or mean-pooling of boundary tokens — fail to model intra-span token interactions, which matters disproportionately in CTI where entities are long, irregular, and semantically composite ("remote code execution"). Its response is a CNN-with-gating span encoder plus orthogonal regularisation of the shared metric space, reaching baseline-beating F1 on 22 fine-grained APTNER classes with 10\% of the training annotations. Beyond the architecture, LETNER contributes an evaluation axis the benchmark literature lacks: a cost-aware framework modeling annotation cost against logarithmic performance growth to locate the point of diminishing returns — making the annotation-budget decision itself an object of study, which is precisely the framing this survey's central thesis predicts the field needs (§1, §8).

The bi-encoder split~\citep{stepanov2026millionlabel} is not a one-off scaling trick — an independent, contemporaneous paper in the adjacent biomedical domain validates the same architectural bet with considerably more statistical rigor. GLiNER-BioMed domain-adapts both GLiNER architectures (uni- and bi-encoder) to biomedicine via a synthetic-data pipeline (distill a large biomedical LLM's annotation behavior into a smaller model, generate high-coverage synthetic pretraining text, then fine-tune on curated real data), and reports every comparison with a Wilcoxon signed-rank test rather than point estimates alone. Its large model reaches 59.77\% zero-shot micro-F1 across eight biomedical benchmarks, +5.96 points over the next-best baseline (p$<$0.001) and +11.22 over the strongest generative-LLM baseline despite far fewer parameters. The bi-encoder result specifically is the more interesting data point for CTI: it beats the uni-encoder at small (+4.40 F1) and base (+3.94 F1) scale but \emph{loses} to it at large scale ($-$4.87 F1), which the authors attribute to doubled encoder capacity mattering only while the uni-encoder is still capacity-starved. Read against the headline throughput numbers of \citet{stepanov2026millionlabel}, this is a caution rather than a contradiction: the bi-encoder's efficiency gains at large label counts and its F1 profile across model scale are separate claims, and a future CTI bi-encoder deployment should not assume the accuracy story transfers just because the throughput story does.

\subsection{LLM-based generative extraction}

Generative extraction inherits the instruction-tuned lineage (InstructUIE, UniNER, GoLLIE — via \citealp{zaratiana2023gliner}) and offers what encoders cannot: free-form task definitions, multi-step reasoning over the text, and output in arbitrary structured formats. The CTI-specific evidence, however, splits sharply by evaluation realism. Mezzi, Massacci, and Tuma~\citep{mezzi2025large} show that the field's optimistic precision numbers (0.83-0.89) were measured on inputs of 18-106 words — sentences shorter than a paper abstract — while real threat reports average $\sim$3,000 words. On 350 full-length APT reports, precision falls to $\sim$0.76 with degradation tracking input length, and few-shot learning and fine-tuning only partially recover it. Worse, the failure modes compound for automation: the same prompt issued twice can extract \emph{different CVEs} (non-determinism that directly corrupts patch prioritization), and confidence is miscalibrated in both directions, so threshold-based filtering cannot rescue the pipeline. The diagnosed mechanism is contextual confusion — real reports mention other APTs, historical campaigns, and alternative vectors that the model cannot exclude as irrelevant.

Two 2025-2026 lines show the unreliability is a property of \emph{naive prompting}, not of LLMs as a method family. AZERG~\citep{lekssays2025text} decomposes text-to-STIX into four fine-tuned subtasks — entity detection, entity typing, related-pair detection, relationship typing — and reaches F1 84-95 on full reports, beating prompted open and closed models by 2-25\%: task decomposition plus task-specific fine-tuning addresses precisely the monolithic-prompt failure the unreliability study measures. Minerva~\citep{alam2026minerva} goes further by exploiting a structural property of the domain: CTI outputs are largely \emph{verifiable} — an ATT\&CK technique ID, a CWE label, a CVSS vector is either canonical or it is not — so reinforcement learning with verifiable rewards applies without learned reward models or preference data. Across 16 verifier-paired training tasks, four open-weight backbones, and 12 benchmarks, MinervaRL gains +15.8 points over base models. The thesis-level observation: the taxonomies of §2.3, built for interoperability, turn out to double as free supervision — a second route (alongside distant supervision, §5.4) by which the field substitutes existing structure for unaffordable annotation.

For deployment, the resulting picture is conditional rather than verdictive. Prompted LLM extraction on full reports remains unreliable in ways that compound silently (length degradation, inconsistency, miscalibration); decomposed fine-tuning and RLVR-trained specialists materially close the gap; and ontology-constrained generation (OntoLogX, §2.3) adds formal validity guarantees at the output boundary. The open question the benchmarks of §8 quantify: even the strongest current models remain weak on the reasoning-intensive end of CTI (attribution, mitigation), independent of extraction accuracy.

A general-domain result quantifies a variable the CTI-specific studies above treat as fixed: output \emph{format}. Zhan, Wang, and Huang compare five ways of asking an LLM to emit NER results — inline-bracketed, inline-XML, and three JSON variants — while holding model and data constant across eight open LLMs on four benchmarks. The gap is not incremental: inline formats average 91.9 F1 (flat/nested combined), while offset-based JSON — asking the model to emit precise character start/end positions — collapses to 53.9/34.5. With the best format, an 8B open model (LLaMA3.1) beats every pretrained encoder baseline on CoNLL2003. This is independent confirmation of a failure mode already documented empirically in our own pipeline (§1): Gemini's \emph{returned} character offsets verify against the source text only $\sim$7\% of the time, which this result reframes as structural rather than implementation-specific — asking a generative model for exact offsets is close to the worst possible extraction contract, regardless of which model is asked. The same paper's low-resource result is a second CTI-relevant data point: LLM performance on the biomedical benchmark GENIA lags the best specialized model by 4 points, the largest gap in the paper, suggesting the length-and-domain-scarcity problems \citet{mezzi2025large} diagnoses for CTI reports are not CTI-specific but a general property of generative extraction under limited domain exposure.

\subsection{Data-centric methods}

Where families 1–6 change the model, data-centric methods change the corpus, and they compose with any model family. Three strategies have emerged, distinguishable by where the new supervision comes from. Harmonization (§3) manufactures scale from \emph{existing labels} by mapping fragmented schemas onto a shared ontology. Description-based augmentation borrows \emph{authoritative definitions}: AnnoCTR shows MITRE ATT\&CK concept descriptions are effective few-shot training augmentation \citep{lange2024annoctr}. Synthetic generation manufactures the \emph{labeled text itself}: SynthCTI~\citep{ruizrdenas2025synthcti} targets the long-tail class imbalance in technique-mapping corpora by clustering each underrepresented class's sentences (HDBSCAN over sentence embeddings), extracting per-cluster semantic features (LDA topics, KeyBERT keyphrases, tone), and prompting a small LLM (Gemma-3 4B) to generate class-faithful synthetic sentences.

SynthCTI's results quantify how much the data side is worth: macro-F1 for ALBERT jumps from 0.35 to 0.52 (+48.6\% relative) and SecureBERT from 0.44 to 0.66 on CTI-to-MITRE — and small augmented models outperform larger unaugmented ones, meaning augmentation bought more than architecture at this data scale. The caveat is task shape: these gains are on sentence classification, where the synthetic example needs only a correct class label. Generating \emph{span-annotated} synthetic text is harder — but no longer open in the adjacent privacy domain: GLiNER2-PII~\citep{zaratiana2026gliner2pii} trained a winning span detector on nothing but constraint-driven synthetic data, where programmatic constraints from the label inventory (entity counts, exclusions) and diversity constraints (locale, format, register) condition LLM generation of text with exact span annotations. The CTI analogue — synthesizing span-annotated threat-report text from a STIX-derived label inventory — has, to our knowledge, not been attempted, and is among the most direct research openings this survey identifies (§9).

A fourth strategy targets a variable orthogonal to all three above: not which labels exist, not what text carries them, but the \emph{composition} of a fixed-size training set. Ye et al. partition CyNER (a human-annotated, not weakly-labeled, CTI corpus) into hard and easy sentences via three interpretable, corpus-adaptive difficulty signals — multi-token entity length, lexical rarity relative to a data-driven threshold, and within-sentence entity-type diversity — cross-checked by an ensemble of the compared models, then retrain at five hard:easy ratios holding total entity count constant. The 1:1 ratio wins on F1 and recall across three backbones and five random seeds (95\% CI widths generally under 0.06); all-easy training inflates precision but measurably degrades recall, meaning a corpus curated toward "clean" examples produces a model that systematically undercounts real entities. The result is a training-time lever any of this section's other strategies can adopt without modification: dictionary/MITRE-derived weak labels already implicitly separate easy sentences (single, common, unambiguous entities) from hard ones (long, rare, mixed-type spans) — the open question this survey flags for §9 is whether a difficulty-aware resampling pass over \emph{weak} labels reproduces this paper's gains, or whether the effect depends on the labels being clean to begin with.

\section{The Adjacent Task: TTP Extraction and ATT\&CK Mapping}

Mapping text to MITRE ATT\&CK tactics and techniques is frequently conflated with CTI NER but is structurally different: techniques are \emph{concepts} that may be mentioned implicitly. AnnoCTR~\citep{lange2024annoctr} makes the distinction operational — an explicit mention names the technique more or less literally, while an implicit one describes the behavior ("send phishing e-mails with malicious attachments" $\rightarrow$ T1566.001 Spearphishing Attachment) with no taggable name span. The corpus quantifies how common the implicit case is: implicit TECHNIQUE mentions outnumber explicit ones roughly 2:1 (2,292 vs 1,098). This is why entity-linking formulations (micro-F1 $\approx$ 65 on technique identification in AnnoCTR) outperform sequence tagging for this task, and why granularity has progressed from document-level labels (rcATT) through sentence-level (TRAM) to span-level with explicit/implicit marking (AnnoCTR). AnnoCTR's linking experiments make the mention/concept distinction quantitative: zero-shot entity disambiguators (BLINK, GENRE) already link GROUP, MALWARE, and TOOL mentions at 82--100\% accuracy, and fine-tuning on the corpus mostly \emph{hurts} those classes (catastrophic forgetting), yet the same zero-shot models collapse to 12--15\% on techniques, where in-domain retraining is essential to reach 66.9\% \citep{lange2024annoctr}. Name-like entities link on surface form; concept-like entities do not, because the textual surface routinely shares nothing with the knowledge-base title.

The task's formulations are best organized by the stratified task space of Kumar et al.~\citep{crossman2026constructing}: eight types from multiclass tactic classification through multi-label hierarchical tagging to text-to-text generation, into which the lineage slots cleanly — TRAM as multi-label technique classification, rcATT as mixed multi-label, TTPDrill as multiclass hierarchical, and the recent LLM work as text-to-text. The same paper supplies the cluster's most sobering empirical result: TF-IDF features with multi-label hierarchical classical models reach $\sim$94\% tactic and $\sim$82\% technique accuracy, meeting or beating state of the art with no LLMs at all — while GPT-4o manages roughly 60\% at the tactic level. Together with the systematic review of 80 studies~\citep{tamanna2026what}, which finds technique-level classification dominant, transformers ascendant, LLM adoption merely emergent, and reproducibility poor (proprietary datasets, unreleased code), the picture is of a subfield where method sophistication and operational performance have not yet aligned.

Three 2025-2026 lines push past the single-sentence, single-report, single-standard framing that the review identifies as the field's default. First, evaluation realism: Haque et al.~\citep{haque2026single} replicate a 29-method comparison at the \emph{campaign} level — 90 reports across SolarWinds, XZ Utils, and Log4j — and find that aggregating reports improves technique F1 by $\sim$26\% with saturation at 5-15 reports, while ceilings stay low (54.9-78.6 F1) and a third of misclassifications confuse semantically similar techniques within the same tactic. Their control-coverage metric converts extraction error into operational exposure: the best method still misses $\sim$23\% of ground-truth mitigations. Notably for corpus builders, technically dense, low-readability reports (MITRE, CISA advisories) contribute most to extraction performance. Second, instruction explicitness: TTPrompt~\citep{peng2025retrieval} demonstrates that retrieval-based in-context learning succeeds largely through incidental entity-type overlap in retrieved examples — implicit induction is unreliable — and replaces it with a hierarchical instruction framework whose levels mirror the TTP abstraction itself, plus feedback-driven refinement that adapts annotation guidelines from errors on $\sim$1\% of training data, approaching full-dataset fine-tunes. Third, standard faithfulness: the evolvable-memory framework~\citep{meng2025instantiating} observes that systems fit \emph{datasets} whose labels drift from the ATT\&CK standard, propagating inconsistency; it converts official definitions into dual-layer situational knowledge — a shared-context retrieval layer and a contrastive layer distinguishing confusable techniques (T1132 vs T1071) — lifting Qwen2.5-32B 18\% Technique-F1 over GPT-4o while making assignments auditable. The contrastive layer directly targets the same-tactic confusion class the campaign study quantifies; the memory's evolvability answers ATT\&CK's own versioning drift.

The cluster's through-lines connect back to the survey's spine. The implicit/explicit-mention distinction (AnnoCTR, §2.1) explains why linking and classification formulations beat span tagging here; the standard-faithfulness critique is the §2.3 taxonomy problem surfacing at the supervision level; and the field's three best levers — aggregation across reports, explicit instructions derived from the standard, contrastive disambiguation — all exploit structure that exists outside any labeled dataset, consistent with the annotation-scarcity thesis. Adjacent surveys frame the positioning: the ATT\&CK applications survey~\citep{jiang2025mitre} covers the framework's uses broadly, the systematic review~\citep{tamanna2026what} covers TTP extraction, and the Büchel SoK (USENIX Security 2025) covers extraction methods comparatively. This survey differs in unifying span NER, TTP mapping, and KG construction under one extraction-pipeline frame, with method-family depth on zero-shot and data-centric approaches.

\section{From Entities to Knowledge Graphs}

Entity extraction is rarely the end goal: the structured artifact CTI pipelines actually want is a knowledge graph — typed entities connected by typed relationships, queryable across reports and time. This section covers the construction path (NER $\rightarrow$ relation extraction $\rightarrow$ triple assembly $\rightarrow$ graph population); graph reasoning, graph-based retrieval, and credibility assessment over completed graphs are out of scope.

The relation extraction step has recently become tractable in the same zero-shot, encoder-based paradigm as modern NER. GLiNER-Relex~\citep{stepanov2026glinerrelex} extends the GLiNER architecture to joint NER and RE in a single forward pass: entity-type and relation-type labels are encoded with the text in one shared context, recognized spans are paired (exhaustively or with adjacency-guided pruning), and pairs are scored against relation-type embeddings using a mechanism the authors trace to knowledge-graph embedding methods. In end-to-end zero-shot evaluation it leads document-level (DocRED, 31.3 micro-F1 vs 18.6 for GPT-5-mini) and cross-domain (CrossRE) benchmarks among comparable systems — significant for CTI, where relations routinely span sentences ("the group deployed X… the malware then exfiltrated via Y"). Its weakness on FewRel's 100 fine-grained relation types is equally instructive: ATT\&CK-flavored CTI relation vocabularies are similarly fine-grained, and no CTI-specific RE benchmark yet exists to measure this.

Construction systems divide into two paradigms whose contrast is cleanest in a single lineage. AttacKG~\citep{li2021attackg} is the pre-LLM reference: pipeline NER extracts attack behavior graphs from reports, technique identification works by aligning extracted graphs against template graphs built from ATT\&CK procedure text, and cross-report aggregation produces technique knowledge graphs — F1 0.887/0.896/0.789 on entities/dependencies/techniques across 1,515 reports, beating Extractor and TTPDrill. ThreatKG~\citep{gao2022threatkg} generalizes this into operational infrastructure: multi-source gathering, per-type extraction models, and continuous graph maintenance — KG construction as a system lifecycle, not a one-shot extraction. The LLM wave then rebuilt both. AttacKG+~\citep{zhang2024attackgboosting} replaces AttacKG's pipeline with four prompted modules (rewriter, parser, identifier, summarizer) and enriches the schema to a temporally unfolding event whose steps each carry a behavior graph, TTP labels, and a state summary. CTINexus~\citep{cheng2024ctinexus}, from the ThreatKG group, replaces both fixed rules and fine-tuning with optimized in-context learning — demonstration retrieval, hierarchical entity alignment, long-distance relation completion — adapting to new ontologies from a handful of examples. A more measured assessment comes from Fieblinger et al.~\citep{fieblinger2024actionable}: open-source LLMs extract triples well under constrained decoding or fine-tuning, but scaling graph completion remains unsolved. MultiKG~\citep{wang2024multikg} widens the input beyond text entirely, merging report-, audit-log-, and static-code-derived graphs into unified technique representations.

The 2026 frontier addresses this survey's oldest theme at the graph level. Every system above defines or upgrades its own schema — AttacKG+'s temporal event model, CTINexus's flexible ontologies, MultiKG's fine-grained technique representation — recapitulating §3's label fragmentation one level up, exactly as hypothesized. ANCHOR~\citep{kim2026schemaagnostic} makes the diagnosis explicit (schema-specific pipelines need manual rebuilds per ontology; prompt-stuffing fails on large ontologies like UCO; enterprise APIs conflict with sensitive incident data) and offers the structural fix: schema-agnostic construction via hybrid ontology discovery — search-and-navigate over arbitrary formal schemas — with SHACL validation enforcing compliance, evaluated across UCO, STIX, and MALOnt, with local models matching enterprise APIs. GRID~\citep{huang2026grid} attacks the supervision side with the same strategy the NER literature adopted (§5.7): generating it synthetically — unsupervised article-graph alignments plus a scripted task bank of cheap verifiable rewards train 4B extractors to 68.53 average F1 across a five-source benchmark, near-top performance at a fraction of the size. Verifiable rewards (Minerva, §5.6), synthetic supervision (§5.7), and small-model economics (§5.5) all converge here.

This section's scope relative to prior reviews: the standing cybersecurity-KG review~\citep{liu2022review} catalogs application scenarios across nine categories with pre-LLM construction coverage; this survey instead follows the text-to-graph construction path through the 2024-2026 LLM wave and connects it backward to the NER and schema literature the KG papers themselves rarely engage. The open problems this section contributes to §9: KG ontology fragmentation (confirmed, with schema-agnosticism as the emerging answer), entity canonicalization across sources (CTINexus's alignment stage meeting Hesperus's naming chaos), and graph-completion scaling~\citep{fieblinger2024actionable}.

\section{Evaluation and Benchmarks}

Evaluation in CTI extraction operates at two levels that the literature often conflates. The first is span-level NER metrics — exact-match precision/recall/F1, strict versus loose boundaries, micro versus macro averaging — where the methodological hazards are cross-dataset ones: Ja\l{}ocha et al.'s 4x4 protocol (§3) remains the only systematic generalization stress test, and its finding that in-dataset F1 wildly overstates portability should be read as a standing caveat on every single-corpus result in this survey. The second level is task-capability benchmarking, which asks not "can the model tag spans" but "can it do the analyst's job," and here the CTIBench lineage (\citealp{alam2024ctibench}; NeurIPS 2024) defines the landscape: knowledge MCQs, CVE-to-CWE root-cause mapping, CVSS severity prediction, ATT\&CK technique extraction, and threat-actor attribution with aliases controlled — a task taxonomy spanning memorization through abductive reasoning.

Its successor AthenaBench~\citep{alam2025athenabench} addresses the failure mode every static benchmark in a fast-moving domain eventually exhibits: staleness and training-data contamination. Tasks are regenerated from live sources — the NVD API and ATT\&CK STIX data with configurable publication windows — so models are evaluated on vulnerabilities and techniques that postdate their training by construction. This is the concrete form of the temporal-validity problem our open-problems ledger flags: in CTI, a benchmark is a snapshot of an adversarial landscape, and its evidentiary value decays accordingly. The headline results temper the LLM optimism of §5.6 from a second direction: across twelve models including GPT-5 and Gemini-2.5 Pro, structured prediction (severity vectors) is largely solved while reasoning-intensive tasks (attribution, risk mitigation) remain subpar even for frontier models, with open-source models trailing badly throughout.

What the benchmark landscape still lacks is instructive. There is no CTI-specific relation extraction benchmark (the GLiNER-Relex gap, §7); no benchmark measures label efficiency or annotation cost (LETNER's PoDR framework stands alone, §5.5); none measures extraction consistency or calibration, the two failure axes the unreliability study (§5.6) showed are decisive for automation; and span-level NER evaluation remains corpus-bound with no AthenaBench-style dynamic refresh. Each gap is an inherited cost of evaluating extraction as a static NLP task rather than as a component of a live intelligence pipeline.

\section{Open Problems}

The fragmentation complex. The field's deepest problem appears at every level of abstraction with the same structure: annotation schemas encode incompatible local ontologies (cross-dataset F1 collapse, \citealp{jalocha2025label}), token- and label-definitions drift between corpora ('sample', 'backdoor', 'Dridex' typed contradictorily), vendors cannot converge on actor names because agreement requires sharing private telemetry~\citep{roa2025hesperus}, KG systems each define their own schema (§7), and the standards themselves are ontologically confused, conflating artifacts with events~\citep{adams2026maecolite}. Partial fixes exist per level — STIX harmonization~\citep{echchammakhy2025cyberner}, entity alignment~\citep{cheng2024ctinexus}, schema-agnostic construction~\citep{kim2026schemaagnostic} — but the root cause is economic: fragmentation is an equilibrium of a market built on proprietary telemetry, so technical normalization can manage it, not eliminate it.

The annotation economy. Expert CTI annotation remains too expensive to scale, and the field's most inventive work is designed to circumvent it: distant supervision from clean security gazetteers (§5.4), zero-shot types-as-language (§5.5), instruction refinement from 1\% of labels~\citep{peng2025retrieval}, verifiable-reward training that converts canonical identifiers into free supervision~\citep{alam2026minerva}, and synthetic generation — proven for classification~\citep{ruizrdenas2025synthcti} and for span annotation in the adjacent PII domain~\citep{zaratiana2026gliner2pii}. Concrete unattempted direction: constraint-driven synthesis of span-annotated CTI text from a STIX-derived label inventory. Meanwhile annotation-budget optimization has exactly one formal treatment (LETNER's point-of-diminishing-returns framework) and no benchmark coverage.

Reliability, robustness, and trust. Prompted LLM extraction degrades with report length, is inconsistent across identical queries, and is miscalibrated in both directions~\citep{mezzi2025large}; no benchmark measures consistency or calibration. Extraction pipelines are adversarially fragile at ingestion — generated fake intelligence passes specialized and LLM filters at 97\%/75\% FPR~\citep{shafee2025false}. Standard-faithfulness is unmeasured: systems fit datasets whose labels drift from ATT\&CK itself~\citep{meng2025instantiating}. And model confidence cannot arbitrate any of this while calibration is unquantified — the same gap CENSOR identified inside distant supervision.

Evaluation realism. Single-report evaluation understates achievable campaign-level coverage by $\sim$26\% and hides the same-tactic confusion class driving a third of errors~\citep{haque2026single}; static benchmarks grow stale and contaminated, with live-source regeneration~\citep{alam2025athenabench} the only current counter; there is no CTI relation-extraction benchmark despite RE being the KG bottleneck (§7); and corpus-bound span-F1 remains the field's default metric despite §3's demonstration that it does not transfer.

Expressiveness gaps. Implicit technique mentions resist span formulations \citep{lange2024annoctr}; multi-label entity roles cannot be expressed by single-label span models (LETNER's own admission); and the source substrate itself carries vendor bias that no downstream method corrects~\citep{suarezroman2026cti}. A small general-domain pilot study (Latifi, 2025 — n=119 tokens, single annotator, no inter-annotator agreement computable, cited here as directional evidence only, not as a rigorous benchmark) is nonetheless a clean illustration of a failure mode CTI inherits structurally: capitalized common nouns and polysemous strings force a type decision context alone must resolve ("Hope" as a name versus an idea in that study; "Apache" as threat actor, tool, or web server in CTI). The one clear signal — LLM context-modeling wins concentrate specifically on person/name disambiguation (0.923 vs 0.692 F1 for traditional taggers) rather than generalizing across entity types (traditional tools still win on ORGANIZATION, 0.733 vs 0.657) — is consistent with SSNER's SCE module (§5.3) treating polysemy as a distinct problem from span detection rather than a byproduct of better embeddings alone. No CTI-specific study isolates this ambiguity class the way this pilot does for general text; doing so at CTI scale (Apache, Cobalt Strike as tool-or-actor, Anonymous as group-or-adjective) remains unaddressed by any paper in this corpus.

Each problem above traces to cited evidence; together they suggest the field's next phase is less about extraction accuracy than about the surrounding system: supervision economics, schema governance, robustness, and evaluation that matches operational reality.

\section{Conclusion}

This survey followed CTI's extraction pipeline from named entities through relations to knowledge graphs, organizing methods by supervision source and type-binding mechanism rather than chronology. Three findings organize the field.

First, fragmentation — of schemas, names, ontologies, and vendor coverage — is the field's persistent structural problem, with the same economic root at every level.

Second, the annotation-scarcity thesis: the field's apparently disparate innovations (harmonization, distant supervision, zero-shot matching, instruction refinement, verifiable rewards, synthetic generation) constitute a single sustained effort to substitute existing structure for unaffordable expert labels — and the structure being exploited is, ironically, the standards infrastructure whose fragmentation causes the other problem.

Third, the encoder-versus-LLM question resolves, on this evidence, into an engineering question of task shape: compact schema-conditioned encoders dominate well-typed extraction at production economics; decomposed, verified, or constraint-bound LLMs earn their cost on implicit mentions, relations, and reasoning-adjacent tasks; and naive prompting is consistently inadequate.

What the field most lacks is not better extractors but better surroundings for them: legally reusable corpora, schema governance with semantics, benchmarks that measure consistency and survive time, and supervision pipelines that treat annotation as the scarce resource it is.

\appendix
\section{Artifact Availability}

Table~\ref{tab:artifacts} indexes the public artifacts of the surveyed
method and dataset papers. ``None'' means no release was found in the
source text; URLs and licenses shown were resolved and checked at
survey time.

\begin{table*}[t]
\centering
\footnotesize
\setlength{\tabcolsep}{4pt}
\begin{tabular}{>{\raggedright\arraybackslash}p{3.7cm}>{\raggedright\arraybackslash}p{2.2cm}>{\raggedright\arraybackslash}p{6.6cm}>{\raggedright\arraybackslash}p{2.0cm}}
\toprule
Work & Type & Artifacts & License \\
\midrule
GLiNER~\citep{zaratiana2023gliner} & NER model & \path{github.com/urchade/GLiNER} + HF models + Pile-NER data & Apache-2.0 \\
GLiNER2~\citep{zaratiana2025gliner2} & multi-task IE & \path{github.com/fastino-ai/GLiNER2}, pip, HF weights & Apache-2.0 \\
GLiNER bi-encoder~\citep{stepanov2026millionlabel} & NER at scale & \path{hf.co/knowledgator/gliner-bi-large-v2.0} (+ edge/small/base); \path{github.com/Knowledgator/GLinker} & Apache-2.0 \\
GLiNER-Relex~\citep{stepanov2026glinerrelex} & joint NER+RE & \path{hf.co/knowledgator/gliner-relex-large-v1.0} (+ base/multi, HF collection) & Apache-2.0 \\
GLiNER2-PII~\citep{zaratiana2026gliner2pii} & PII spans & \path{hf.co/fastino/gliner2-privacy-filter-PII-multi} & open \\
GLiGuard~\citep{zaratiana2026gliguard} & moderation & \path{github.com/fastino-ai/GLiGuard} & open \\
LETNER (ESWA 2026) & CTI NER & none — reimplementation from equations required & — \\
BOND~\citep{liang2020bond} & DS-NER & \path{github.com/cliang1453/BOND} (code + distant labels) & open \\
CENSOR~\citep{si2023improving} & DS-NER & \path{github.com/PKUnlp-icler/CENSOR} & open \\
TTPrompt~\citep{peng2025retrieval} & prompted NER & none found & — \\
AnnoCTR~\citep{lange2024annoctr} & corpus & \path{github.com/boschresearch/anno-ctr-lrec-coling-2024} & CC-BY-SA 4.0 \\
CyberNER~\citep{echchammakhy2025cyberner} & corpus & \path{github.com/yasirech-chammakhy/CyberNER} & public \\
AttackER~\citep{deka2024attacker} & corpus & \path{zenodo.org/records/10276922} (data); \path{hf.co/Cyber-ThreaD} (fine-tuned models only) & CC-BY-4.0 (data) \\
AZERG~\citep{lekssays2025text} & STIX extraction & \path{hf.co/datasets/QCRI/AZERG-Dataset} (+ T1--T4 Mistral models, QCRI collection) & Apache-2.0 \\
Label unification~\citep{jalocha2025label} & eval harness & \path{github.com/PLtier/NLP-Cyber-NER} & unverified \\
SynthCTI~\citep{ruizrdenas2025synthcti} & augmentation & no code release; builds on \path{github.com/dessertlab/cti-to-mitre-with-nlp} (upstream dataset); reimplementable & CC-BY-SA-4.0 (upstream data) \\
Minerva~\citep{alam2026minerva} & RLVR & no public release located as of July 2026 & — \\
CTIBench~\citep{alam2024ctibench} & benchmark & \path{github.com/xashru/cti-bench} & open \\
AthenaBench~\citep{alam2025athenabench} & benchmark & \path{github.com/Athena-Software-Group/athenabench} & open \\
JPMC tagger~\citep{crossman2026constructing} & TTP tagging & \path{github.com/jpmorganchase/MITRE_models} & public \\
Evolvable memory~\citep{meng2025instantiating} & TTP extraction & none found & — \\
AttacKG~\citep{li2021attackg} & KG construction & \path{github.com/li-zhenyuan/Knowledge-enhanced-Attack-Graph} & MIT \\
ThreatKG~\citep{gao2022threatkg} & KG system & none confirmed & — \\
CTINexus~\citep{cheng2024ctinexus} & KG construction & \path{github.com/peng-gao-lab/CTINexus} & open \\
AttacKG+~\citep{zhang2024attackgboosting} & KG construction & "upon acceptance" — not confirmed & — \\
MultiKG~\citep{wang2024multikg} & KG construction & none confirmed & — \\
Actionable CTI~\citep{fieblinger2024actionable} & KG construction & none confirmed & — \\
ANCHOR~\citep{kim2026schemaagnostic} & KG construction & none confirmed & — \\
GRID~\citep{huang2026grid} & KG construction & not confirmed (preprint) & — \\
HiP / Hesperus~\citep{roa2025hesperus} & actor naming & \path{github.com/0xjet/hip} & open \\
OntoLogX~\citep{cotti2025enabling} & ontology+LLM & none found & — \\
Unreliability study~\citep{mezzi2025large} & evaluation & none (data via Di Tizio et al., open) & — \\
\bottomrule
\end{tabular}
\caption{Reproducibility index for method and dataset papers.
``None'' = no release found in the source text.}
\label{tab:artifacts}
\end{table*}


\paragraph{Citation note.} External works referenced only through the
surveyed papers (e.g., SecureBERT, UniNER, rcATT, TRAM) are cited as
secondary sources; full reading notes for the 41-paper corpus are in the
repository's \path{notes/bibliography.md}.

\bibliographystyle{acl_natbib}
\bibliography{survey}

\end{document}
